% ===== §7 =====
\section{The Fourth Impossibility: The Constitution of the Subject}
\label{p25:sec:constitution}

\noindent
The three prior impossibilities concerned the objective: whether it can be written, fixed from without, and furnished with a common scale. This last one reaches beneath the objective to the quietest condition of the anatomy, the one that underwrites the others and is least often brought to light, that there is an \emph{optimizer}. Maximization requires not only something to be maximized but someone to maximize it: a subject whose objective $J$ is, whose preferences it encodes, standing outside the problem, prior to it and stable through it, as the one for whom $x^{\ast}$ is best. Without such a subject the phrase ``his optimum'' has no determinate referent, for there is no settled \emph{his} to whom the optimum belongs. This section denies that condition. In intimate relation the subject is not prior to the relation but constituted within it; the one who would optimize is a product of the process he is imagined to stand outside of; and the ``who'' the whole construction presupposes loses the fixity the construction needs.

\subsection*{The subject as constituted, not given}

The optimizing frame inherits from marginalism a subject who arrives at the situation already formed, bearing a preference order settled in advance, which the situation then reveals and satisfies. Even the more sophisticated economics that allows preferences to change tends to model the change as itself the output of a prior and stable meta-preference, a fixed disposition governing how tastes evolve, so that the fixity merely retreats one level and the subject remains, at bottom, given before the process. The relational ontology of the series denies this picture at its root, and denies it of the subject exactly as it denied it of value, since the two denials are one. The subject of intimate relation does not enter the relation with its preferences fixed and undergo the relation as their playing-out. It is formed in the relating: its desires are shaped by the other's desire, its sense of what is worth wanting is remade in the shared life, and its individuation, the very drawing of the boundary between what it is and what the other is, is accomplished and revised within the bond rather than settled before it \parencite{paperxxiv,paperxvii}. This is not the mild and assimilable point that people change their minds, which a meta-preference could absorb. It is the ontological point that the preference-bearing subject the optimization presupposes, the fixed evaluator standing over against the field of options, is not the antecedent condition of the intimate relation but one of its ongoing products. What the relation generates is not only value between the partners but the partners themselves, as the valuers they are in the course of becoming.

\begin{casebox}{who would collect the optimum}
A person enters a relation with a settled sense of what he wants from a life, and over years within it that sense is remade: he comes to want things he had not known were wantable, to find hollow some things he had once ranked highest, to care about the other's flourishing in a way that is not reducible to his own prior preferences because it has reorganized them. Suppose now he had, at the outset, tried to choose the relation by optimization, computing which arrangement best satisfied the preferences he then held. The arrangement that maximized the preferences of the man at the outset is not the arrangement the man he became would choose, and not because either made an error, but because the second man's preferences are the product of the very relating that the first man's optimization was meant to govern from before it began. Whose optimum, then, was to be computed? Not the first man's, for he no longer exists to collect it, and his preferences were dissolved and reformed by the process. Not the second man's, for he did not exist yet at the moment of choice, and could not have supplied the objective, being himself an output of what was chosen. There is no single subject present throughout to be the owner of the objective and the beneficiary of its maximum. There is a subject in the course of being constituted by the relation whose value the optimization would compute, which is a different thing from the fixed optimizer the frame requires.
\end{casebox}

\subsection*{The dissolution of the fixed referent}

From this a difficulty follows that is fatal to the fourth condition, and it is a difficulty about reference rather than about stability, which the case has made ready to state. If the optimizer is constituted within the very process his optimization is supposed to govern, then ``whose optimum?'' has no fixed answer, because the subject named in any answer differs at the end of the process from the one who posed it at the start, and the difference is produced by the process itself. This is worse than the drift of the objective treated under the second impossibility, and must not be folded into it. There the target moved while the archer stood still; here the archer is remade by the shooting, so that the one who would gather the optimum at the summit is not the one who set out to climb and did not hold, at setting out, the preferences by which the summit is a summit. The frame needs a stable subject to be the owner of the objective and the beneficiary of its maximization, the one whose $J$ it is and for whom $x^{\ast}$ is best; and in intimate relation there is no such stable owner, only a subject in the making whom the relation is precisely in the business of unfixing. To ask what is optimal \emph{for him} is to presuppose a settled him that the relating dissolves and reforms, and the referent of the optimizing ``who'' comes apart in the asking.

\begin{claim}[The constituted subject]
Optimization presupposes an \emph{optimizer}: a subject prior to the problem and stable through it, whose objective the function is and for whom its maximum is best. In intimate relation the subject is constituted within the relation rather than given before it, its preferences and its individuation ongoing products of the relating. The optimizer is therefore a product of the process it is imagined to stand outside of, and ``whose optimum?'' loses a fixed referent, the subject named in any answer being remade by the very process the optimization would govern. The fourth condition of the anatomy fails: there is no stable subject to own the objective or to be the beneficiary of its maximum.
\end{claim}

\subsection*{The four taken together}

The four impossibilities may now be gathered, and gathering them shows that they are not four names for one objection but four independent failures at four distinct joints, running from the surface of the problem to its ground. Can the values be \emph{written} into the objective? Not all: the drive, extinguished by attainment, admits no term. Can the objective, so far as written, be \emph{fixed} from without? No: relational value admits no exterior standpoint, and the objective is endogenous of ontological necessity. Can the values it fixes be \emph{compared} on a common scale? Not the incommensurable ones: the scale $\argmax$ needs is not there. And is there a stable \emph{subject} to own the whole and receive its maximum? No: the optimizer is constituted in the process it would govern. Each cut is at a different joint, and each alone suffices, since an optimization problem needs all four conditions and the failure of any one leaves the operation undefined. Together they establish, without a single appeal to value, that intimate relation furnishes no well-defined optimization problem: not a hard one that greater computational power would crack, but an ill-posed non-problem, whose required object, a representationally complete, exogenous, fixed, scalar objective belonging to a stable subject, the operation of maximization does not find. With this the formal movement of the paper is complete. What follows is no longer proof but criticism, and it turns to the different question the proof has cleared the way for: not whether one \emph{can} optimize intimate relation, which has been answered in the negative, but what befalls the relation when, undeterred by the answer, one imposes the frame regardless.
